\documentclass[12pt]{article}
\usepackage[utf8]{inputenc}
\usepackage[T1]{fontenc}
\usepackage{hyperref}
\title{Wadjiginy Language Resources for K-12 Curriculum Development}
\author{Compiled by Claude Sonnet 5 (claude-sonnet-5) for Andr\'as Kornai}
\date{}
\begin{document}
\maketitle

\section{Introduction}
Wadjiginy (ISO 639-3 wdj; Glottocode wadj1254) is also called Wagaydy, Wogait or Batjamalh. It is a moribund Australian Aboriginal language of the Northern Territory, from the Anson Bay area near the Daly River mouth, north of Darwin. It is closely related to the extinct Kandjerramalh (Pungupungu) and together they form the Wagaydyic group. One source gives 5 speakers (2005). The search results did not link it to Wagaman.

\section{Primary Educational Resources}
No school primers or programmes were found in search results.

\section{Linguistic and Reference Materials}
The main reference named in search results is Lysbeth Ford's \emph{Batjamalh Grammar and Dictionary} (2016), cited by the AIATSIS Austlang entry: \url{https://collection.aiatsis.gov.au/austlang/language/N31}. Overview pages: \url{https://en.wikipedia.org/wiki/Wadjiginy_language} and \url{https://en.wikipedia.org/wiki/Wagaydyic_languages}. I did not see the grammar itself and cannot confirm where it is available.

\section{Religious and Community Resources}
None found in search results.

\section{Digital and Audio Resources}
None found in search results. Archival recordings may exist with AIATSIS, but this was not confirmed.

\section{Curriculum Development Context}
The grammar and dictionary are the key foundation for any revitalization work. With so few speakers, community consultation with Wadjiginy people is essential.
\end{document}